(a) For $P\in Z$, the intersections $V\cap Z$, with $V$ ranging over neighborhoods of $P$ in $X$, are all neighborhoods of $P$ in $Z$. Thus $(i_*\mathscr F)_P=\mathscr F_P$. For $P\notin Z$, there is a neighborhood disjoint from $Z$, so the stalk is zero.

(b) If $P\in U$, its neighborhoods contained in $U$ are cofinal, so the defining presheaf has stalk $\mathscr F_P$. If $P\notin U$, every neighborhood has value zero. Sheafification preserves these stalks and restricts to $\mathscr F$ on $U$.

If $\mathscr G$ is another sheaf with the stated restriction and stalks, map the defining presheaf to $\mathscr G$ by the chosen identification on opens contained in $U$ and by the zero map on other opens. Sheafification gives $j_!\mathscr F\to\mathscr G$, which is an isomorphism on every stalk and hence an isomorphism. It is unique with the chosen restriction, since its stalk maps are already determined everywhere.

(c) The first map is induced by the identity on opens contained in $U$, and the second is the natural restriction morphism $\mathscr F\to i_*i^{-1}\mathscr F$. At $P\in U$, their stalk sequence is $0\to\mathscr F_P\to\mathscr F_P\to0\to0$; at $P\in Z$, it is $0\to0\to\mathscr F_P\to\mathscr F_P\to0$. Both are exact, so the sequence of sheaves is exact by \exref{II.1.2}.
